\begin{claim}{F.7}{Open}{The nonsolvable centralizer comparison and its checked finite image}
For an oriented mapping class group $M$ of genus at most two,
the normalizer and centralizer comparisons for finite solvable
subgroups have the complete proof below. For a geometric
$A_5\le M=\Gamma^+_{0,\nu}$, the equality
\[
                 C_{\widehat M}(A_5)
                    =\overline{C_M(A_5)}                 \tag{F.7.1}
\]
is stated in \cite[Theorem~C(iii)]{R-Boggi} and
\cite[Theorem~5.1(iii)]{R-BZ}. The displayed proofs contain
an additional closure--intersection equality for which we
have not verified a justification. We do not assert that
(F.7.1) is false.

There is nevertheless equality of their images in the
finite puncture-permutation quotient. Write the invariant
puncture set as a union of some of the exceptional orbits
of sizes $12,20,30$ and $m$ free orbits of size $60$.
Both images fix the exceptional punctures individually
and induce the full group $A_5^m\rtimes S_m$ on the
free orbits. In particular, with only exceptional orbits
the unresolved part of (F.7.1) lies in the pure profinite
mapping class group; its discrete right side is trivial.

F.12--F.14 identify the residual twelve-puncture $A_5$ centralizer with an exact common-involution equalizer in the pure eight-point completion and prove injectivity after the specified ordinary-point forgetting. F.15 proves the full two-point pushing comparison and removes every cocycle in the actual intersection of the two forgetting kernels. Its correcting elements form the explicitly described fixed-subgroup cosets. The original source image has trivial intersection with that kernel; membership of a corrected element in that image remains the full equalizer equation. F.16 constructs a full closed procyclic subgroup of the correcting centralizer that every proper labelled forgetting kills and whose intersection with the source image is trivial. The exact membership test using all proper forgetting images, even with the common-involution centralization equation, therefore has a strictly larger solution group than the original source image. F.17 identifies the sixteen additional punctures required by the literal rational fixed-point drilling and proves its nonfactorization through the original source, using primitive full procyclic winding detectors. F.18 proves that every original residual element has a unique full branch-ordered centralizing lift. Its membership in the rational drilling image is the distinguished fibre of the explicit continuous coset map $\mathfrak o:\mathcal E\to\mathcal J/K^+$. Whether that fibre is the whole residual, and whether the residual is trivial, remain unresolved.
F.19 identifies the joint forgetting kernel $D^+$ as the closed ambient normal closure of eight actual cross-collision meridians and proves its exact one-class cocycle comparison with a unique corrector. The canonical correction therefore lies in $D^+$, and the source-membership obstruction has the precise target $D^+/K^+$. The actual drilled image and the full two-track pullback have common explicit splittings over the original rank-eleven source, with kernels $K^+$ and $D^+$ respectively. Equality of those two subgroups and triviality of the original residual remain unresolved.
F.20 proves the actual rank-three cross-divisor coset coordinate, its eleven based source actions and all original transports. Exact configuration computations and finite group calculations prove transitivity on the $82$ invariant $A_5$ quotients, all $3,321$ distinct paired quotients, and six noncentral quotients of order $60\cdot2^{121}$. The specified continuous finite wreath maps consequently have the same full images on $K^+$ and $D^+$, for every coefficient prime. The full source-membership equation remains unresolved.
F.21 proves transitivity on all eight full cross-divisor coset spaces and constructs lifts of their entire divisor groups in the actual selected-meridian centralizers. Consequently the closed normal closure of $K^+$ inside $D^+$ is $D^+$, and $K^+$ maps onto the full abelianization and maximal pronilpotent quotient of $D^+$. The exact equality $K^+=D^+$ and the original residual source-membership equation remain unresolved.
F.22 gives the actual two-point kernel as a semidirect product of two explicit compactly generated profinite kernels, retains its full Artin action, and identifies the interaction kernel killed by both one-point forgettings. The source map has image exactly $K^+$; its joint one-point image is the product of its two actual marginal images. Fullness of those marginals and surjectivity of the displayed source map remain unresolved.
F.23 gives an explicit based source-sign symmetry carrying the two actual marginal inclusions isomorphically to one another. It specifies all twenty-seven original source images by exact nine-strand words, with the original basepoint paths and inverse fiber coordinates. For every continuous homomorphism from either named ambient rank-seven group to $S_3$ or $A_4$, the marginal source and marginal target have equal images. The full marginal equalities, the interaction source image and the original residual membership remain unresolved.
F.24 proves injectivity of the actual restriction on continuous rank-one crossed homomorphisms and first cohomology for every finite-field character nontrivial on the marginal target. Sixty-four actual source rows have rank seven on four explicit Laurent charts over every characteristic. The signed comparison retains the original source maps. Every scalar affine representation has the same fixed-point set on the marginal source and target, and its two marginal images agree whenever the source image meets the translation subgroup trivially. The remaining nonzero translation intersections, full marginal images and interaction source image remain unresolved.
F.25 proves equality of those affine images for every representation whose source translation intersection is linear over the ambient scalar field. This includes every representation with equal ambient and marginal scalar fields, and every representation into $\operatorname{Aff}_1(\mathbb F_p)$ for every prime $p$. The chosen scalar lift, its translation conjugation, and the separating coefficient maps are explicit. Unequal scalar fields without that linearity property, the full marginal equalities and the interaction source image remain unresolved.
F.26 gives an exact equivalent trace-cocycle kernel criterion for a proper affine marginal image, with the actual conjugation and coefficient homomorphism. Its tensor and Frobenius coefficient maps, including their inverses, identify the remaining rank target as $6m+1$ on $7m$ columns. Complete finite verification and an independently written replay prove equality of the actual marginal images for every scalar field of order nine, even inside a larger original coefficient field. General attainment of the trace rank, the full marginal equalities and the interaction source image remain unresolved.
F.27 proves the exact trace rank and affine marginal image equality over every finite coefficient field on the actual two-parameter character locus $(1,s,t,1,1,1,(st)^{-1})$. Its integer Laurent commutators force an essential source line, and actual character-one seed conjugates give an invertible Frobenius Vandermonde comparison. F.28 extends this to the three-parameter locus $(r,s,t,1,1,1,(rst)^{-1})$: an explicit two-dimensional invariant inclusion and four quotient coordinates separate all minimal source submodules across distinct Frobenius factors. The individual restriction theorem then proves summed restriction injective and rank $6m+1$. Both results retain the exact based sign map and include larger original coefficient fields. The general scalar-character locus, full marginal equalities, interaction source image and original residual membership remain unresolved.
F.29 proves the larger exact four-parameter locus
$(r,s,t,u,1,1,(rstu)^{-1})$ over every finite coefficient field.
Four actual principal-coordinate charts retain all original source words.
Three charts have complete comparisons of their minimal invariant
subspaces across Frobenius factors; the exceptional overlap has a
forced-lift residual with coordinate $-2$ in its proved odd
characteristic. The fourth chart has an essential line detected by
actual character-one seed conjugates and an explicit Vandermonde
inverse. Summed restriction has rank $6m+1$, and the exact trace
and based sign maps give both marginal image equalities on this locus.
The general character locus and the original full marginal and
interaction comparisons remain unresolved.
\end{claim}
\begin{proof}[Source audit and proved comparisons]
The supplied preprint snapshot of \cite{R-Boggi}
has since appeared in \emph{S\~ao Paulo Journal of
Mathematical Sciences}, volume~20, article~35,
published 23 September 2026, DOI
\texttt{10.1007/s40863-026-00558-y}.
The publisher's full text was checked as well:
Proposition~2.4 retains the same displayed
closure--intersection step, and Lemma~3.2
still invokes that proposition.
The audit below applies to both inspected versions;
it is a proof-step audit and makes no claim based
on the paper's publication status.

For every finite subgroup $J\le M$, there are canonical
continuous maps
\[
 \widehat{C_M(J)}\longrightarrow
       \overline{C_M(J)}\lhook\joinrel\longrightarrow
       C_{\widehat M}(J).                              \tag{F.7.2}
\]
The first is onto by density and compactness. Its
injectivity states that the topology induced from $M$
on the discrete centralizer is its full profinite
topology. Surjectivity of the second arrow is the
ambient centralizer comparison. These are two specific
properties of the displayed maps; injectivity of the
first arrow does not supply surjectivity of the second.

In \cite[Proposition~2.4]{R-Boggi}, choose the displayed
generators $x_1,\ldots,x_k$ of the finite group, each of
prime-power order. The prime-power centralizer theorem
does prove
\[
 C_{\widehat M}(J)
   =\bigcap_i C_{\widehat M}(x_i)
   =\bigcap_i\overline{C_M(x_i)}.                      \tag{F.7.3}
\]
The source then identifies this last intersection with
$\overline{\bigcap_i C_M(x_i)}$. The same final equality
occurs in the proof of
\cite[Theorem~5.1(iii)]{R-BZ}, there indexed by all
$x\in J$. Neither displayed argument supplies a
closure--intersection hypothesis. This exact step is
what is left unverified here. The later virtually
compact special application in \cite[\S6.2]{R-BZ}
has a separate relatively quasiconvex product-separability
input; no such input has been proved here for the
punctured-sphere mapping class group.

\emph{A noncircular solvable proof.}
We first record the prime-power inputs used in the
induction. Goodness and finite virtual cohomological
type of mapping class groups of genus at most two,
and their finite-index subgroups, permit the direct
application of \cite[Theorem~A]{R-BZ}. Thus for a
finite $p$-subgroup $A$ we have
\[
 N_{\widehat M}(A)=\overline{N_M(A)},\qquad
 C_{\widehat M}(A)=\overline{C_M(A)}.                   \tag{F.7.4}
\]
The proof of \cite[Theorem~A and Corollary~B]{R-Boggi}
also supplies the full induced profinite topology on
$C_M(A)$ and hence on $N_M(A)$ in this prime-power case.
Here is why this use avoids the preceding gap. Its
free-action argument is the restriction-injectivity
argument on the finite covering surface group. Its
ramification-removal argument uses Proposition~2.4
only for the specified group $A$, now a $p$-group,
in the good configuration-group extension. For that
application the second equality in (F.7.4), directly
from \cite[Theorem~A]{R-BZ}, replaces Proposition~2.4's
intersection argument. Thus the two required inputs
are justified without using any nonsolvable assertion.
The passage from centralizer to normalizer topology
uses its finite-index kernel under conjugation into
$\Aut(A)$; this kernel is open in the induced topology,
as is detected by any finite quotient in which $A$
injects.

We prove the solvable normalizer comparison, together
with its full induced topology, by induction on the
order of the finite solvable group, simultaneously for
all finite-index subgroups of these mapping class
groups. Finite-index passage preserves the comparison:
intersect the completed normalizer with the corresponding
open subgroup and use density in an open set.

For a nontrivial finite solvable $J$, take a nontrivial
$p$-primary subgroup $A$ of the last nonzero term of
its derived series. That term is abelian, so its
$p$-primary subgroup is characteristic in $J$.
Consequently every normalizer of $J$ normalizes $A$.
The discrete covering normalizer sequence is
\[
 1\longrightarrow A\longrightarrow N_M(A)
                 \longrightarrow Q\longrightarrow1,   \tag{F.7.5}
\]
where $Q$ is a finite-index subgroup of the mapping
class group of the quotient surface with its branch
and puncture decorations. Its genus is again at most
two. The exact quotient map is induced by descent of
the commuting covering diagram; the decorations
record the original punctures, the removed branch
points and the covering subgroup. This is the discrete
normalizer sequence in \cite[\S2.1]{R-Boggi}.

By the prime-power topology input, the completed
normalizer in (F.7.4) is $\widehat{N_M(A)}$.
Completing (F.7.5) therefore gives its quotient by
the finite embedded kernel $A$ as $\widehat Q$.
The subgroup $J/A\le Q$ is smaller and solvable.
The induction hypothesis applies to its normalizer
in $\widehat Q$. The inverse image of that normalizer
under the quotient in (F.7.5), on either side, is
exactly the normalizer of $J$: $J$ is the full inverse
image of $J/A$ because it contains $A$. Pulling back
the dense discrete normalizer proves
\[
                  N_{\widehat M}(J)=\overline{N_M(J)}.
                                                               \tag{F.7.6}
\]
It also proves the topology assertion. Indeed the
completion of this inverse image maps to
$\widehat{N_M(A)}$; on the quotient by $A$ it is
injective by induction, and on $A$ it is the specified
inclusion. Its kernel is therefore trivial.
This completes the induction.

Conjugation gives a continuous homomorphism
\[
       N_{\widehat M}(J)\longrightarrow\Aut(J)
\]
to a finite group. Its kernel is the ambient centralizer
and is open. The dense subgroup $N_M(J)$ from (F.7.6)
meets that open kernel densely, and its intersection
there is exactly $C_M(J)$. This proves
\[
                  C_{\widehat M}(J)=\overline{C_M(J)}
\]
for every finite solvable $J$, without taking an
unjustified intersection of closures.

\emph{The finite image for $A_5$.}
Realize the geometric $A_5$ as the icosahedral rotation
group. Its exceptional point stabilizers have orders
$2,3,5$, and the corresponding transitive orbits are
$A_5/C_2$, $A_5/C_3$, $A_5/C_5$. Every permutation
commuting with this transitive action is right
multiplication by an element of $N_{A_5}(C_d)/C_d$.
This quotient has order two in each of the three
cases; its nonidentity permutation exchanges the
two poles of each relevant rotation axis.

An orientation-preserving homeomorphism commuting
with an odd-order nonidentity rotation fixes its
two poles individually. The local oriented rotation
numbers at the two poles are opposite, and an
orientation-preserving commuting homeomorphism
preserves this local rotation number. This excludes
the nonidentity permutation on the $12$- and
$20$-point orbits for the discrete cyclic centralizers.
Their completed centralizers have the same finite
permutation images by (F.7.4). The full $A_5$
centralizer lies in these cyclic centralizers, so
it fixes these exceptional punctures individually.

For an order-two stabilizer choose a tetrahedral
subgroup $A_4\le A_5$ containing it. An
orientation-preserving homeomorphism commuting
with $A_4$ fixes every ramification point of the
$A_4$ quotient. To see this, its quotient fixes the
branch point of order two. It also fixes each of the
two branch points of order three, rather than
exchanging them, by the preceding local rotation
argument. Isotope the quotient map to the identity
relative to these three branch points and lift the
isotopy. The lift stays in the centralizer; its
endpoint is a deck transformation in $Z(A_4)=1$.
The ramification set is finite and invariant along
the isotopy, so every one of its points is fixed.
The usual discrete covering-centralizer comparison
provides a commuting representative for a centralizing
mapping class. Since $A_4$ is solvable, the proved
centralizer comparison transfers this fixed-puncture
property to $C_{\widehat M}(A_4)$. It therefore holds
for $C_{\widehat M}(A_5)$, excluding the last possible
nonidentity permutation on the $30$-point orbit.

On $m$ free orbits the finite commutant is
$A_5^m\rtimes S_m$. Every element of this commutant
is already realized by the discrete centralizer.
For a direct construction, remove the three branch
points from the quotient sphere. Its unramified
$A_5$-cover has surjective monodromy onto $A_5$.
Move one of the $m$ quotient punctures around a
loop, keeping all the others fixed, and lift the
isotopy starting at the identity. The lift commutes
with $A_5$ and acts on that one free orbit by its
specified monodromy, while fixing the other orbits
and all ramification points. These loops realize
each factor $A_5$. Half-twists interchanging the
quotient punctures realize $S_m$ modulo these factors.
Thus the displayed wreath product is the exact
finite permutation image on both sides of (F.7.1).

The program
\path{code/part_a/a5_centralizer_permutations.py}
constructs $A_5$ as the sixty even permutations of
five letters. It checks its order-two and order-three
generating pair, the three coset actions, normalizer
orders $4,6,10$, and commutants of order two. The
fixed-pole constraints leave only the identity in
each exceptional block; the order-two block uses
the six ramification points for the specified $A_4$.

For completeness, the residual intersection question
has an exact finite-quotient formulation. Take
$u=(1\ 2)(3\ 4)$ and $v=(1\ 3\ 5)$ in $A_5$, which
generate it, and put $D_u=C_M(u)$, $D_v=C_M(v)$.
For open normal subgroups $V\subset U$ of $M$, let
$D_u(V),D_v(V)$ denote their images in $M/V$.
Compactness of coset fibres proves
\[
 \operatorname{im}\bigl(C_{\widehat M}(A_5)\to M/U\bigr)
 =\bigcap_{V\subset U}
       \operatorname{im}\bigl(D_u(V)\cap D_v(V)\to M/U\bigr).
                                                               \tag{F.7.7}
\]
Indeed choose a lift $t\in\widehat M$ of a point
in the right side. For each $V\subset U$, the set
\[
 (\overline{D_u}\,\overline V)\cap
 (\overline{D_v}\,\overline V)\cap t\overline U
\]
is nonempty and compact. These sets are directed
under refinement, so their finite intersections
are nonempty. Compactness produces a single element
in their intersection, which is
$\overline{D_u}\cap\overline{D_v}\cap t\overline U$.
Conversely any such element supplies every finite
coordinate. Equations (F.7.3) and (F.7.7) identify
the remaining assertion precisely: its right side
must equal the image of $C_M(A_5)$ for every $U$.
The finite permutation computation proves that
equality at the puncture-permutation level. It does
not establish it at every refinement. With only
exceptional orbits the discrete centralizer is
trivial, since the quotient sphere has just the
three differently ordered branch points and
$Z(A_5)=1$; the still unproved assertion is then
that the pure completed centralizer is trivial.

\emph{The minimal twelve-puncture case.}
There is a further exact localization of this remaining group.
Write $P=P\Gamma_{0,12}$ and choose $A=A_4\leq A_5$.
The action of $A$ on the twelve vertices is free and transitive.
Its quotient has one free puncture and three branch points of
orders $2,3,3$. The labelled commuting homeomorphism group on
the filled sphere is identified, by descent and the unique
commuting lift, with the identity homeomorphism component of
the quotient fixing those three branch points. It is
contractible. Evaluating the free quotient puncture consequently
identifies the mapping-class centralizer with the free group
$\pi_1(S^2\setminus\{3\text{ branch points}\})=F_2$.
Its permutation on the free fibre is the surjective monodromy
onto $A_4$. The pure centralizer is therefore the exact kernel
\[
 1\longrightarrow C_P(A)=F_{13}\longrightarrow F_2
                    \longrightarrow A_4\longrightarrow1.
                                                               \tag{F.7.8}
\]
Indeed it is the fundamental group of
$X=S^2\setminus R_A$, where the $A_4$ ramification set $R_A$
has $6+4+4=14$ points; its rank is also
$1+12(2-1)=13$ by Schreier. Label the vertices by $g\in A_4$
using a chosen vertex $z_0$. The inclusion is the configuration
map $z\mapsto(gz)_{g\in A_4}$, followed by the quotient by
M\"obius transformations. The proved solvable comparison and
full topology identify its completion with $C_{\widehat P}(A)$.
In particular the unknown group $C_{\widehat M}(A_5)$ lies in
this specified embedded $\widehat F_{13}$.

Each of the six subgroups $C_5\leq A_5$ has two fixed vertices
and two free five-cycles. Retain either one fixed vertex and
one five-cycle, giving twenty-four six-point sets, or retain
both fixed vertices and one five-cycle, giving twelve
seven-point sets. For each retained set $I$, the continuous
forgetting map $r_I$ is defined on $\widehat P$ and is equivariant
for the specified $C_5$. In either target the quotient has
exactly its two branch points and one free puncture. Its
discrete $C_5$ centralizer is therefore exactly the deck group
$C_5$: the quotient mapping class group fixing these three
decorated points is trivial. Prime-power comparison gives the
same completed centralizer. Its intersection with the pure
group is trivial, since the deck group permutes the five-cycle.
It follows that
\[
 C_{\widehat M}(A_5)\subset
 C_{\widehat P}(A)\cap
       \bigcap_{I\text{ among these }36\text{ sets}}\ker(r_I).
                                                               \tag{F.7.9}
\]

The six-point maps in (F.7.9) have a proved common kernel even
on the discrete $F_{13}$. The antipodal permutation on the
twelve vertices commutes with $A_5$. In the regular $A_4$
labelling it is right multiplication by an involution $h\in A_4$.
Every retained six-point set contains exactly one member of
each antipodal pair. Collisions of two entries $g_i z,g_j z$
at a ramification point with stabilizer $H\leq A_4$ occur
precisely when $g_j^{-1}g_i\in H$; these are the right cosets
$gH$. Thus there is no collision among the retained entries
at either fixed point of $h$. Each six-point configuration map
extends across both of these deleted points of $X$. Their
nontrivial peripheral meridians in $F_{13}$ map to the identity
under all twenty-four maps. Completion retains these meridians
as nontrivial primitive procyclic subgroups. Consequently the
joint six-point map on the embedded $\widehat F_{13}$ is
noninjective by these explicit elements.

Every seven-point set retains a collision at each of the
fourteen ramification points: seven choices from six pairs
force a pair, and seven choices from four triples force a
triple collision. The program verifies all twelve sets and
the preceding common invisible axis for the six-point sets,
using right cosets throughout. It also computes the induced
maps on $H_1(F_{13},\F_{101})$. A meridian at an $A_4$
ramification point maps to the sum of the pair generators
within its collision clusters. In the target pure sphere
group the pair generators satisfy the vertex-sum and total-sum
relations. The two fixed points of each cyclic stabilizer give
the same homology column. Each seven-point map and their joint
map have rank six, with seven-dimensional kernel; the joint
six-point map has rank five. These are exact calculations of
the indicated abelian quotients. They neither produce a
nontrivial $A_5$ centralizer element nor prove injectivity of
the joint seven-point map on $\widehat F_{13}$.

F.8 proves that all twelve seven-point kernels in (F.7.9)
are the same, both discretely and after completion. It further
proves that the discrete kernel contains no nonidentity reducible
mapping class and that the completed kernel meets every conjugate
of every full source cusp subgroup trivially. The exact four-point
characters detecting those cusps and all eighty-four finite
subset--inertia checks are supplied there. F.9 subsequently
constructs an explicit nontrivial common kernel element with
a full procyclic closure, and proves that all twenty-four
six-point restrictions factor through these seven-point maps.
The thirty-six-fold kernel is therefore nontrivial. The full
centralizer comparison and geometric conjugacy remain unproved;
no converse to the inclusion (F.7.9) is inferred.

\medskip\noindent\textbf{The exact common-$C_3$ comparison.}

The following maps give an exact model for the residual group in (F.7.1).
The fixed intersection displayed below is the remaining unresolved group.
All permutations use the five-letter identification of $A_5$ already
fixed in F.7--F.10.  Put
\[
 A=\operatorname{Stab}_{A_5}(5),\qquad
 r=(1\ 2\ 3),\qquad K=\langle r\rangle,
 \qquad\beta=(1\ 2)(4\ 5).
\]
Then $A\cong A_4$, $\beta K\beta^{-1}=K$ by inversion,
and $\langle A,\beta\rangle=A_5$.

\medskip\noindent\textbf{The actual quotient homomorphism.}
Let $P=P\Gamma_{0,12}$ at the icosahedral configuration.
The discrete branched-cover sequence for the $K$ action is
\[
 1\longrightarrow K\longrightarrow N(K)
       \longrightarrow Q\longrightarrow1,
\]
where $Q$ is the finite-index liftable subgroup of the
mapping class group of the quotient sphere with two decorated
branch points and four free punctures.  The centralizer
$C_P(K)$ has finite index in $N(K)$: both the centralizing
condition and the pure puncture condition have finite index.
Its intersection with $K$ is trivial, since nonidentity deck
maps permute the twelve punctures.  Its image $J$ in $Q$ is
therefore isomorphic to $C_P(K)$ and has finite index in $Q$.
Every element of $J$ fixes the four free quotient punctures.
It also fixes the two branch points separately: exchanging
them conjugates the order-three local deck generator to its
inverse, and hence cannot centralize it.  Thus
\[
             C_P(K)\xrightarrow{\ \sim\ }J
                         \leq P\Gamma_{0,6},
               \qquad [P\Gamma_{0,6}:J]<\infty.          \tag{F.7.10}
\]
The prime-power comparison and full induced topology proved
in F.7 identify $C_{\widehat P}(K)$ with
$\widehat{C_P(K)}$.  Completing the finite-index embedding
(F.7.10) identifies it with the specified open subgroup
$\widehat J\leq\widehat{P\Gamma}_{0,6}$.  This gives an
injective continuous homomorphism on the actual completed
centralizer; no completed intersection identity is used.

By F.7.8 the orbit map identifies $C_P(A)$ with the specified
$F_{13}=\pi_1(X)$, where
\[
 X=\mathbf P^1\setminus R_A,\qquad
 R_A=\{\infty\}\cup\{z:z(z^4-1)(z^8+14z^4+1)=0\}.
\]
Since $K\leq A$, composition with (F.7.10) is the injection
$j:F_{13}\hookrightarrow P\Gamma_{0,6}$ described next.
The solvable comparison and full topology for $A$ and $K$
show that its continuous extension
\[
                \widehat j:\widehat F_{13}
                       \hookrightarrow\widehat{P\Gamma}_{0,6}
                                                               \tag{F.7.11}
\]
is injective.  This is a different, explicitly specified map
from the thirty-six forgetful maps whose common kernel was
computed in F.9.

\medskip\noindent\textbf{The six labelled quotient points.}
In the F.9 tetrahedral coordinate,
\[
 r(z)=\frac{z+i}{z-i},\qquad
 q(z)=z+r(z)+r^2(z)
     =\frac{z^3+3iz+1-i}{(z-1)(z-i)}.
\]
The degree-three quotient $q$ is invariant under $r$.
Its branch points upstairs are
\[
 p_\pm=\frac{(1+i)(1\pm\sqrt3)}2,
\]
and their images are $3p_\pm$.  In the twelve-point
labelling of F.8 the four free $K$-orbits, in order, are
\[
 (0,8,1),\quad(11,3,5),\quad(2,6,9),\quad(7,10,4).
                                                               \tag{F.7.12}
\]
Their representatives have coordinates $z,1/z,-1/z,-z$.
Consequently the unnormalized labelled quotient map is
\[
 z\longmapsto
   (3p_+,3p_-,q(z),q(1/z),q(-1/z),q(-z)).                \tag{F.7.13}
\]
Its entries are distinct on $X$: any equality between two
free quotient points would identify their points in the
$A$ orbit, and a collision with a branch point would give a
nontrivial stabilizer in $A$.

Put $U(z)=(z-p_+)/(z-p_-)$.
Direct substitution gives
\[
       \frac{q(z)-3p_+}{q(z)-3p_-}=U(z)^3.
\]
Normalizing the first, second, and third points in (F.7.13)
to $0,\infty,1$ gives the actual map
\[
 \Sigma:X\longrightarrow M_{0,6},\qquad
 z\longmapsto(s_1(z),s_2(z),s_3(z)),
\]
where
\[
 s_1=\left(\frac{U(1/z)}{U(z)}\right)^3,
 \qquad
 s_2=\left(\frac{U(-1/z)}{U(z)}\right)^3,
 \qquad
 s_3=\left(\frac{U(-z)}{U(z)}\right)^3.                 \tag{F.7.14}
\]
The homomorphism $j$ in (F.7.11) is precisely $\Sigma_*$,
by descent of the same labelled orbit path.

\medskip\noindent\textbf{The second $A_4$ and the exact based involution.}
The involution $\beta$ exchanges the two branch points and
permutes the four free orbits in (F.7.12) by $(0\ 1)(2\ 3)$.
On the normalized moduli coordinates its relabelling is
therefore
\[
 B(s_1,s_2,s_3)=(s_1,s_1/s_3,s_1/s_2).                 \tag{F.7.15}
\]
For example, after swapping the branch points the new
normalizing coordinate is $u\mapsto s_1/u$, since the
first free point has become the old point $s_1$.
This derives (F.7.15) from the actual labelled quotient,
rather than selecting a conjugate linear action.

Use $z_*$ and the path
$\gamma:z_*\to5+z_*\to b=5+4i$ from (F.10.12).
The point $\Sigma(z_*)$ is fixed by $B$.  Let
$\tau=\Sigma\circ\gamma$, with the indicated direction.
For a chronological loop $\ell$ based at $\Sigma(b)$, the
transported involution is exactly
\[
 \beta_b(\ell)=
 \tau^{-1}*B(\tau)*B(\ell)*B(\tau)^{-1}*\tau.          \tag{F.7.16}
\]
The path endpoints in this expression match in order;
$B^2=1$ and $B\Sigma(z_*)=\Sigma(z_*)$ show directly that
$\beta_b^2=1$.  It induces the conjugation action of the
same transported finite mapping class $\beta$ used in F.10
on the injective quotient image (F.7.10).

Since $A$ and $\beta$ generate $A_5$, the exact residual
comparison is
\[
 \widehat j\bigl(C_{\widehat P}(A_5)\bigr)
   =\widehat j(\widehat F_{13})\cap
       \operatorname{Fix}\!\left(
          \beta_b:\widehat{P\Gamma}_{0,6}
                    \longrightarrow\widehat{P\Gamma}_{0,6}
                          \right).                    \tag{F.7.17}
\]
Indeed centralizing $A_5$ is exactly centralizing $A$ and
$\beta$.  For elements centralizing $A$, equation (F.7.11)
is injective; also $\beta$ normalizes $K$, so both the
element and its $\beta$ conjugate lie in the domain of
that injection.  Equality of their quotient images is
therefore equivalent to their equality upstairs.

\medskip\noindent\textbf{Executed exact checks and their scope.}
The program \path{code/part_a/second_a4_c3_attack.py}
uses exact arithmetic in $\Q(i,\sqrt5,\sqrt{\varphi+2})$
to check the $60$ M\"obius matrices, $180$ generator
products, and $720$ vertex-label images.  It verifies all
four orbits in (F.7.12), all representative functions,
the quotient identity, and (F.7.14).  It also checks
the fixed-locus calculation
\[
 \Sigma(z)\in\operatorname{Fix}(B)
 \quad\Longleftrightarrow\quad s_1(z)=s_2(z)s_3(z),
\]
whose numerator, up to a nonzero scalar and after
cancelling the displayed rational denominators, is
\[
 z(z^2+1)(z^8+16z^6+46z^4+16z^2+1).                    \tag{F.7.18}
\]
It verifies exactly that $z_*$ satisfies the relevant
quartic factor.  Finite intersection of these geometric
loci does not by itself calculate the intersection of
their completed fundamental groups in (F.7.17).

The auxiliary numerical braid discovery in
\path{work/part_a/f9_c3_johnson_probe.py} currently gives
no exclusion at the first nonzero outer Johnson degree.
It is explicitly exploratory: the crossings there have
not been certified for the actual based map (F.7.16),
and a higher-degree comparison must retain that path.
No conclusion about (F.7.17) uses that calculation.

\medskip\noindent\textbf{The modular orbifold and the exact $GT_0^1$ image.}
There is a further exact description of the ambient peripheral
centralizer in the minimal twelve-puncture case.  It also gives
five specified three-cusp tests for the remaining geometric
intersection.  The intersection itself remains unresolved.

Let
\[
 \Delta=\langle a,b\mid a^2=b^3=1\rangle,
 \qquad E=\widehat\Delta=C_2\amalg C_3,
 \qquad d=ab,\quad e=ba.                                \tag{F.7.19}
\]
The isomorphism with the full profinite coproduct follows from
the same universal maps on the two finite factors.
Use the epimorphisms
\[
 \begin{array}{c|cc|c}
     &a&b&\text{kernel}\\ \hline
 q_3:E\to S_3 &(2\ 3)&(1\ 2\ 3)&K\\
 q_5:E\to A_5 &(1\ 2)(3\ 4)&(1\ 3\ 5)&H .
 \end{array}                                             \tag{F.7.20}
\]
The discrete modular orbifold has one order-two point,
one order-three point, and one puncture, whose oriented
loop is $d$.  In the first quotient $d$ has order two,
and in the second it has order five.  The corresponding
torsion-free covers have the following exact data:
\[
 \begin{array}{c|c|c|c|c}
  &\text{degree}&\text{cusp width}&\text{number of cusps}
        &\text{surface type}\\ \hline
  K&6&2&3&(0,3)\\
  H&60&5&12&(0,12).
 \end{array}                                             \tag{F.7.21}
\]
Indeed the quotient maps inject both finite orbifold
stabilizers, so the covers are surfaces.  Their Euler
characteristics are the degree times $-1/6$; the cusp
counts are the indices of $\langle q_i(d)\rangle$ in
the finite quotient.  The resulting genera are zero.
Completion of these finite-index inclusions retains their
full profinite topology.  Thus $K\cong\widehat F_2$,
$H\cong\widehat F_{11}$, and the second extension is
the actual centre-free orbifold extension of the
icosahedral twelve-puncture action.

Here is the precise identification with the specified
action in F.7.  Start with that oriented sphere $S$
and its twelve-vertex puncture set $V$.  Choose a
basepoint $s\in S\setminus V$ outside every nontrivial
stabilizer and use its image in the quotient orbifold.
Choose ordinary positive orbifold loops $u_2,u_3$
around the order-two and order-three points, with
whiskers so that the positive puncture loop is
$u_c=(u_2u_3)^{-1}$.  Set
\[
             a=u_3^{-1}u_2u_3,\qquad b=u_3^{-1};
                   \qquad ab=u_c.
\]
Path lifting gives the monodromy map to the actual
deck group $A_5$.  Its generating pair can be labelled
as the pair $q_5(a),q_5(b)$ in (F.7.20), by the exact
single-orbit calculation below.  This relabels the
deck group and does not change its kernel or the
given subgroup of the sphere mapping class group.
Based loop lifting now identifies the discrete
kernel with $\pi_1(S\setminus V,s)$, and its completion
identifies $H$ with the same specified profinite
surface group.  The chosen lift of the cusp whisker
selects a vertex $v_0\in V$ stabilized by $q_5(d)$.
The actual identification of the twelve labels is
\[
        A_5/\langle q_5(d)\rangle\longrightarrow V,
                   \qquad g\langle q_5(d)\rangle
                                    \longmapsto g v_0.
\]
Under it the cusp $[t d^5t^{-1}]_H$ is exactly the
puncture $q_5(t)v_0$.  In particular the identification
of $H$ carries its pure geometric mapping class
subgroup to the one used in F.7.
Changing the base path changes this identification
by an inner surface-group map; changing the chosen
lift changes it by the specified deck action.
An equivalent choice of the oriented orbifold
presentation gives the corresponding based covering
isomorphism.  These maps conjugate the peripheral
outer group and its geometric subgroup together,
and carry the displayed $A_5$ to the same geometric
action.  They therefore carry the intersection
below to the identical residual centralizer under
that covering identification.

For precision, $x=d^2$ and $y=e^2$ are a full free basis
of $K$.  Schreier rewriting with transversal
$1,a,b,ab,ba,b^{-1}$ gives the following kernel words
$t_g s t_{q_3(gs)}^{-1}$:
\[
 \begin{array}{c|cc}
 t_g&s=a&s=b\\ \hline
 1&1&1\\ a&1&1\\ b&1&1\\
 ab&x&y^{-1}\\ ba&1&y\\ b^{-1}&x^{-1}&1 .
 \end{array}                                             \tag{F.7.22}
\]
Rewriting $a^2$ and $b^3$ at these six representatives
gives only freely trivial words in $x,y$; these twelve
relator checks are included in the certificate.
The resulting Schreier presentation is exactly
$\langle x,y\mid\ \rangle$, giving the stated basis
and then its full profinite basis.  Conjugation by $a$ sends
$(x,y)$ to $(y,x)$, while conjugation by $b$ sends it
to $(y,x^{-1}y^{-1})$.  Writing
$\omega(x)=y$, $\omega(y)=(xy)^{-1}$, the latter
automorphism is exactly $\operatorname{Inn}(y)\omega$.
Therefore the outer
deck action on $K$ is exactly the geometric $S_3$
used in D.35, with these peripheral generators.

\emph{Characteristic kernels and the finite normalization.}
There are six generating order-two/order-three pairs in
$S_3$, forming one orbit under its six inner automorphisms.
There are $120$ such generating pairs in $A_5$, forming
one orbit under conjugation by $S_5$.
For the latter count, fix $(1\ 2)(3\ 4)$: the generating
three-cycles use the fifth letter and one letter from
each of its two pairs, giving eight choices.
Their products with the involution have order five.
The remaining choices generate proper subgroups.
The finite certificate below verifies all $300$ pairs,
as well as the transitive $S_5$ action.
Consequently each epimorphism in (F.7.20) has the unique
kernel of its indicated quotient type.  In particular
both $K$ and $H$ are characteristic in $E$.

We also need the exact quotient-automorphism statement.
The five $V_4$ subgroups of $A_5$ form an intrinsic set.
Its conjugation action is faithful.  An automorphism
acting trivially on these five subgroups acts trivially
on that faithful permutation representation, and hence
on $A_5$.  Thus $\operatorname{Aut}(A_5)$ embeds in $S_5$;
all $120$ conjugations by $S_5$ already give automorphisms,
so this is the full automorphism group.  An odd
conjugation exchanges the two $A_5$ conjugacy classes
of five-cycles.  It follows that an automorphism
preserving the conjugacy class of $q_5(d)$ is inner.
The certificate verifies the five $V_4$ groups,
faithfulness, and both twelve-element five-cycle classes.

\emph{The full peripheral centralizers.}
For $T=K,H$, let $\mathcal P_T^1$ denote the subgroup
of $\operatorname{Out}(T)$ preserving each individual
oriented peripheral element up to conjugacy with
multiplier one.  Put
\[
 \Omega_E^1=\{[\eta]\in\operatorname{Out}(E):
                         \eta(d)\sim_E d\}.
\]
There are canonical continuous isomorphisms, with
the normalizations described below,
\[
 C_{\mathcal P_H^1}(A_5)
       \xrightarrow{\ \sim\ }\Omega_E^1
       \xrightarrow{\ \sim\ }
 C_{\mathcal P_K^1}(S_3)=GT_0^1,
 \qquad GT_0^1=\ker(\lambda:GT_0\to\widehat\Z^\times).
                                                               \tag{F.7.23}
\]
We prove the maps rather than identify the two actions
only by their surface types.

First, for $m=2$ and $T=K$, or for $m=5$ and $T=H$,
\[
                         C_E(d^m)=\overline{\langle d\rangle}
                                      \cong\widehat\Z . 
                                                               \tag{F.7.24}
\]
The element $d^m$ is a primitive cusp of the free group
$T$, so $C_T(d^m)=\overline{\langle d^m\rangle}$.
This follows, for example, by writing $T$ as the
free profinite product of that primitive cyclic factor
and its free complement and using factor malnormality
as in D.35.  The quotient action on the cusp classes
is the actual action on
$Q/\langle q(d)\rangle$, where $Q=S_3$ or $A_5$.
An element centralizing $d^m$ therefore has quotient
in $\langle q(d)\rangle$.  Multiplying it by a suitable
integral power of $d$ puts it in $T$ without changing
centralization.  This proves the equality in (F.7.24).
The procyclic group is a full $\widehat\Z$ because
$d^m$ already generates such a subgroup: a relation
$d^u=1$ would imply $(d^m)^u=1$ and hence $u=0$.

The centralizer $C_E(T)$ is trivial.  Its image in
$Q$ would act trivially in $\operatorname{Out}(T)$,
but that action is faithful on the cusp permutation
set.  Thus this image is trivial, and the remaining
element lies in $Z(T)=1$.
Consequently conjugation embeds $E$ in
$\operatorname{Aut}(T)$ as the full inverse image
of the specified $Q\leq\operatorname{Out}(T)$.

Let $[\phi]\in C_{\mathcal P_T^1}(Q)$.
Conjugation by an actual representative $\phi$ on
$\operatorname{Aut}(T)$ preserves that inverse image
and gives an automorphism $\eta$ of $E$ inducing
the identity on $Q$ and restricting to $\phi$ on $T$.
Changing $\phi$ by an inner automorphism of $T$
changes $\eta$ by the same inner automorphism of $E$.
Also $\eta(d)^m$ is conjugate in $T$ to $d^m$.
After that conjugation, (F.7.24) places $\eta(d)$
in $\overline{\langle d\rangle}$; the equation
$m u=m$ in $\widehat\Z$ forces its exponent to be
one.  Thus $[\eta]\in\Omega_E^1$.

Conversely, choose $[\eta]\in\Omega_E^1$.
Its induced automorphism of $Q$ is inner: for
$S_3$ this follows from its full automorphism group;
for $A_5$ it follows from preservation of the
five-cycle class proved above.
If that quotient automorphism is $\operatorname{Inn}(q_0)$,
choose $u\in E$ with $q(u)=q_0^{-1}$ and replace
$\eta$ by $\eta_0=\operatorname{Inn}(u)\eta$.
The new automorphism induces the identity on $Q$.
Different choices of $u$ differ by an element of $T$,
so the outer class of its restriction to $T$ is
well defined.  Changing the representative of
$[\eta]$ has the same inner effect after normalization;
the centre of $Q$ is trivial, so there is no further
quotient ambiguity.

Write $\eta_0(d)=rdr^{-1}$.  Because the quotient
action is now the identity, $q(r)$ centralizes
$q(d)$.  In both finite groups that centralizer is
exactly $\langle q(d)\rangle$.  Hence $r=t d^j$
for $t\in T$, and $\eta_0(d)=tdt^{-1}$.
For every $s\in E$, the elements $\eta_0(s)$ and
$s$ have the same quotient.  Therefore the restriction
preserves each $T$-cusp class
$[s d^m s^{-1}]_T$ with exponent one.
It belongs to $\mathcal P_T^1$ and centralizes the
outer $Q$ action.

These constructions are inverse homomorphisms.
For injectivity, a normalized automorphism whose
restriction to $T$ is inner can first be adjusted
to fix $T$ pointwise.  Then
$\eta(s)s^{-1}\in C_E(T)=1$ for every $s\in E$,
so it is the identity.  In the other direction,
an inner automorphism of $E$ arising from a
centralizing outer class of $T$ has quotient in
$Z(Q)=1$, hence its conjugator lies in $T$.
Continuity follows from conjugation on the specified
closed subgroup of $\operatorname{Aut}(T)$ and
normalization by the finite quotient.  The maps
are homeomorphisms because all the groups here
are compact Hausdorff.
Finally the equality with $GT_0^1$ is precisely the
Harbater--Schneps centralizer identification used
in D.35, now with the exact basis and outer actions
in (F.7.22)\cite[Propositions 3--4 and Theorem 2]{R-HS}.

Let $\iota:GT_0^1\hookrightarrow\mathcal P_H^1$
denote the inverse of the left-to-right isomorphism
in (F.7.23).  Genus-zero congruence identifies the
completed pure geometric group with its actual
image in $\mathcal P_H^1$.  Thus the remaining group
has the exact description
\[
             C_{\widehat{P\Gamma}_{0,12}}(A_5)
                 =\widehat{P\Gamma}_{0,12}
                                \cap\iota(GT_0^1)
                 \quad\text{inside }\mathcal P_H^1.     \tag{F.7.25}
\]
This proves the relationship with the full $GT_0$
problem; it does not calculate this intersection.

\emph{Explicit substitutions and three-cusp quotient maps.}
Write $\gamma=(1,f)\in GT_0^1$ in its normalized
representative on $K$, so
$x\mapsto x$, $y\mapsto f^{-1}yf$ and $f\in K'$.
Its actual extension to (F.7.19) has the particularly
simple formulas
\[
          \eta_f(a)=f^{-1}a,\qquad
          \eta_f(b)=f^{-1}b,\qquad \eta_f(d)=d .         \tag{F.7.26}
\]
Indeed the first symmetry is $\theta(f)=f^{-1}$.
Direct conjugation on $K$ gives
$\gamma\theta=\operatorname{Inn}(f^{-1})\theta\gamma$,
so $C_E(K)=1$ forces the stated image of $a$.
Since $\gamma(x)=x$, (F.7.24) with $m=2$ forces
$\eta_f(d)=d$; the displayed image of $b=a^{-1}d$
then follows.  These formulas describe the actual
automorphism furnished by (F.7.23), including its
relations and its invertibility.

The quotient automorphism on $A_5$ fixes $q_5(d)$
and is inner.  There is consequently a unique
$j\in\{0,1,2,3,4\}$ for which it is
$\operatorname{Inn}(q_5(d)^j)$; equivalently
\[
             q_5(f)=q_5(e)^j q_5(d)^{-j}.
\]
The pure normalized representative on $H$ is
\[
       \eta_f^H=
          \left.\operatorname{Inn}(d^{-j})\eta_f\right|_H.
                                                               \tag{F.7.27}
\]
For example, with $c_0=d^5$ and $c_1=a d^5a^{-1}=e^5$,
put $h_1=d^{-j}f^{-1}e^j\in H$.  Then
\[
                 c_0\longmapsto c_0,
        \qquad c_1\longmapsto h_1c_1h_1^{-1}.           \tag{F.7.28}
\]
This uses the finite quotient to ensure $h_1\in H$;
it is not a formal substitution of a word outside
the source of a filling homomorphism.

Choose the cusp labels as the twelve left cosets
$A_5/\langle q_5(d)\rangle$, in the breadth-first
order specified by the certificate.  The corresponding
representatives $t_i\in\Delta$ and peripheral words
$c_i=t_i d^5t_i^{-1}$ are recorded there; $t_0=1$
and $t_1=a$.
For a three-element label set $J$ define the actual
continuous filling homomorphism
\[
 q_J:H\longrightarrow H_J
     =H/\langle\!\langle c_i:i\notin J\rangle\!\rangle_H
                         \cong\widehat F_2.            \tag{F.7.29}
\]
The isomorphism type follows from filling precisely
those nine punctures of the specified sphere;
completion preserves the displayed quotient
presentation.  Since (F.7.27) is pure peripheral,
it descends to $H_J$.  Thus there is a continuous
homomorphism
\[
 \mathcal D_J:GT_0^1\longrightarrow\mathcal P_{H_J}^1,
 \qquad\mathcal D_J(\gamma)=
 \bigl[q_J(h)\longmapsto q_J(\eta_f^H(h))\bigr].       \tag{F.7.30}
\]
The brackets on the right specify the outer
automorphism induced by the displayed substitution.
Any change of representative in (F.7.27) changes
it only by an inner automorphism of $H_J$.
Formula (F.7.26), the finite determination of $j$,
and (F.7.29) define each map without assuming that
any two arbitrarily based retained cusp words
form a full free basis.

The pure geometric mapping class group of the
three-punctured sphere is trivial.  The actual
forgetting homomorphism on pure mapping classes,
its action on the filled fundamental group, and
continuity on completions therefore give
\[
 \iota^{-1}\!\left(
       C_{\widehat{P\Gamma}_{0,12}}(A_5)\right)
               \subseteq\bigcap_{|J|=3}\ker\mathcal D_J.
                                                               \tag{F.7.31}
\]
For $g\in A_5$, a lift in $E$ induces the actual
quotient isomorphism $H_J\to H_{gJ}$ by conjugation;
different lifts differ by an inner automorphism.
Because $\iota(\gamma)$ centralizes $A_5$ in
$\operatorname{Out}(H)$, these isomorphisms intertwine
$\mathcal D_J(\gamma)$ and $\mathcal D_{gJ}(\gamma)$.
Their kernels are consequently equal within each
actual $A_5$ orbit of three-element subsets.

The $220$ triples form exactly five orbits, with
the following representatives in the recorded coset
labelling:
\[
 \begin{array}{c|ccccc}
 J&\{0,1,2\}&\{0,1,4\}&\{0,1,8\}&\{0,1,10\}&\{0,4,9\}\\
 |A_5J|&20&60&60&60&20 .
 \end{array}                                             \tag{F.7.32}
\]
The certificate supplies a labelled $A_5$ transporter,
as a word in $a,b$, for every triple.  Hence the
intersection of $220$ kernels in (F.7.31) equals the
intersection of these five specified kernels.
No injectivity or triviality of that intersection
has been proved here.

The standard-library executable
\path{code/part_a/modular_a5_gt0_f7.py} has been run.
It checks the six and $120$ epimorphism pairs and
their single postcomposition orbits, both faithful
cusp actions, both cyclic cusp centralizers, the
five-$V_4$ automorphism argument, the exact twelve
Schreier rewritings in (F.7.22), all twelve peripheral
words, and all $220$ subset transporters in (F.7.32).
These finite calculations certify the finite inputs
and the displayed words.  The centre-free extension,
normalization, completion, and filling maps are
proved above.


\medskip\noindent\textbf{Exact correspondences for the two orbit-twenty detectors.}
The two orbit-twenty tests in (F.7.32) admit the following
explicit subgroup and filling description. It gives
continuous word substitutions and a common section of
the surface-group quotient maps. Faithfulness of the
outer-automorphism detectors remains unresolved.

Use the precise $E,q_5,H,K,a,b,d,e$ of (F.7.19)--(F.7.22).
Put $T=\langle q_5(b)\rangle$ and $B=q_5^{-1}(T)$.
The index-twenty orbifold cover has two order-three
points and four width-five cusps, and no order-two
points. This is the exact coset action: $a$ has ten
two-cycles, $b$ has two fixed cosets and six three-cycles,
and $d$ has four five-cycles. Its signature is
$(0;3,3;4\text{ cusps})$.

The following words give a specified full profinite
presentation:
\[
 \begin{split}
 z&=aba b^{-1},\qquad \tau=zaz^{-1},\\
 r&=b,\qquad s=\tau b\tau^{-1},\\
 u&=(ab^{-1})^4a,\qquad v=(ab)u(ba),\\
 w&=(aba b^{-1}a)(ba)^3(b^{-1}aba).
 \end{split}
 \qquad
 B=\langle r,s,u,v,w\mid r^3=s^3=1\rangle_{\rm prof}.
                                                               \tag{F.7.33}
\]
Thus $B=C_3\amalg C_3\amalg\widehat F_3$ with the
displayed inclusion into $E$. The exact Schreier
certificate uses twenty right cosets $Tg$ and the
prefix-closed breadth-first transversal in $a,b$.
Its forty directed generator edges have nineteen
tree edges, leaving twenty-one Schreier generators.
Rewriting $a^2,b^3$ at the twenty vertices gives
forty relators. Sixteen Tietze eliminations, each
solving a relator in a generator occurring exactly
once, leave the five displayed generators and
precisely the two cube relators.
The complete transversal, original generator words,
substitutions, and eliminations are recorded.
Reduction in $C_2*C_3$ checks the compact words in
(F.7.33). These are inverse presentation maps;
completion of the finite-index discrete inclusion
retains its full profinite topology.

The four based cusp loops and their actual
$T$-orbits of $H$ punctures are
\[
 \begin{array}{c|c|c}
 i&\kappa_i\text{ in }B&J_i\\ \hline
 0&u^{-1}r&\{0,1,2\}\\
 1&v^{-1}&\{3,4,6\}\\
 2&wu&\{5,7,8\}\\
 3&vw^{-1}s&\{9,10,11\}
 \end{array}                                                \tag{F.7.34}
\]
Each is the rewritten word $t_i d^5t_i^{-1}$
for its recorded transversal representative.
The first and fourth triples lie in the
$A_5$-orbit of $\{0,1,2\}$; the middle two lie
in the orbit of $\{0,4,9\}$.

Fill the other three orbifold cusps and put
$B_i=B/\langle\!\langle\kappa_j:j\ne i\rangle\!\rangle_B$.
To
$A=\langle R,S\mid R^3=S^3=1\rangle_{\rm prof}=C_3\amalg C_3$
send $r$ to $R$, $s$ to $S$, and use the indicated row:
\[
 \begin{array}{c|ccc|c}
 i&u&v&w&\kappa_i\\ \hline
 0&S^{-1}&1&S&SR\\
 1&R&S^{-1}R^{-1}&R^{-1}&RS\\
 2&R&1&S&SR\\
 3&R&1&R^{-1}&RS
 \end{array}                                                \tag{F.7.35}
\]
Substitution kills the three required cusp relators.
Conversely those three relators solve for $u,v,w$
as the entries of the row. Inclusion of $r,s$ in
the quotient gives the inverse map from $A$.
This proves the exact quotient isomorphism and
identifies the retained positive cusp as the
displayed product without a power ambiguity.

Let $E^+=\ker(\epsilon:E\to C_2)$, where
$\epsilon(a)=1,\epsilon(b)=0$.
The two-coset Schreier calculation with transversal
$1,a$ gives
\[
 A\xrightarrow{\ \sim\ }E^+,\qquad
                  R\longmapsto b,\quad S\longmapsto aba.
                                                               \tag{F.7.36}
\]
Write $p_i:B\to E^+$ for the composite of the row
(F.7.35) and (F.7.36).
The quotient $B/H=T$ sends
\[
 (r,s,u,v,w)\longmapsto
       (\bar b,\bar b^{-1},\bar b,1,\bar b^{-1}),
\]
whereas $E^+/K=C_3$ sends $b$ to its generator
and $aba$ to the inverse. Identify these two cyclic
groups by their named generators.
All four rows intertwine the specified quotient maps,
and hence $p_i(H)=K$.

The restriction to $H$ is precisely the filling
of the nine punctures outside $J_i$.
Indeed the closed $B$-normal closure of the other
three cusp words is contained in $H$.
Conjugation by $B/H=T$ supplies the three $H$ cusp
classes in each omitted orbit, and further
$H$ conjugation gives their closed normal closure.
Thus
\[
 \begin{split}
 \langle\!\langle\kappa_j:j\ne i\rangle\!\rangle_B
 &=\langle\!\langle c_\ell:\ell\notin J_i
                                      \rangle\!\rangle_H,\\
 1\longrightarrow H_{J_i}&\longrightarrow B_i
                 \longrightarrow C_3\longrightarrow1,
       \qquad H_{J_i}\xrightarrow[\ p_i\ ]{\ \sim\ }K .
 \end{split}                                                \tag{F.7.37}
\]
Each omitted cusp has width one for $H\leq B$:
$\kappa_j\in H$ and its three lifts have precisely
the labels in (F.7.34). No peripheral power is
inserted into this normal-closure identity.

There is a common actual section. Define
\[
 \iota_0:E\longrightarrow E,\qquad
                a\longmapsto\tau,\quad b\longmapsto b .
                                                               \tag{F.7.38}
\]
The orders give a homomorphism. On
$E^+=\langle b,aba\rangle$ it maps the two free
profinite factors to $r,s$, whose inclusion in
$B$ is injective by its retraction onto those factors.
Moreover $q_5(\tau)$ is an involution inverting
$q_5(b)$, so $q_5\iota_0(E)$ is the order-six
normalizer of $T$. Its quotient by $T$ is
exactly $\epsilon$. Thus $\ker\iota_0$ is contained
in $E^+$, where it is trivial, proving that
(F.7.38) is a continuous embedding.

The kernel of
$E\to q_5\iota_0(E)\cong S_3$ is the characteristic
$K$ proved in (F.7.20). Hence $\iota_0(K)\leq H$.
Every $p_i$ is the identity on the displayed two
factors after (F.7.36), so
\[
 p_i\iota_0|_{E^+}=1_{E^+},\qquad
             p_i|_H\circ\iota_0|_K=1_K.                 \tag{F.7.39}
\]
The program also rewrites $\iota_0(x),\iota_0(y)$
through all four quotient rows and verifies that
their images are exactly $x,y$.

For $\gamma=(1,f)\in GT_0^1$, the actual normalized
automorphism (F.7.27) extends to an automorphism of
$E$ inducing the identity on $A_5$. It preserves
$B$ and each filled normal subgroup.
Using (F.7.37) for the target identification, the
actual detector is
\[
 \mathcal D_{J_i}(\gamma)
        =\bigl[p_i\,\eta_f^H\,\iota_0|_K\bigr]
                          \quad\hbox{in }\mathcal P_K^1.
                                                               \tag{F.7.40}
\]
Here $\eta_f^H$ is applied to $\iota_0(K)\leq H$.
Its preservation of that embedded subgroup is
not assumed. The subsequent quotient map is
part of the composition. The composition is
an automorphism because it is the actual descended
automorphism under (F.7.37), evaluated on the
section (F.7.39).
The two pairs of equal detector kernels are
the actual orbit transports of (F.7.31)--(F.7.32).

The four source quotient maps themselves have
a nontrivial common kernel:
\[
 [\kappa_0,\kappa_1]
           =u^{-1}r v^{-1}r^{-1}uv\ne1.                 \tag{F.7.41}
\]
Every row kills this word.
Its nontriviality in the full completion is
detected by $B\to S_3$ sending $r,s,w$ to $1$,
$u$ to $(1\ 2)$, and $v$ to $(2\ 3)$:
its image has order three.
This source word gives an inner automorphism
of $E$, whose class in $\operatorname{Out}(E)$
is the identity. It therefore supplies no
nonidentity $GT_0^1$ element in the kernel of
(F.7.40). Faithfulness of an outer detector,
or of the five-detector family, remains open.

The executed program
\path{code/part_a/modular_c3_correspondence_f7.py}
records the full Schreier presentation, its
sixteen eliminations, the four three-step filling
eliminations, compact word identities, exact
finite quotient labels, eight common-section
word checks, and the finite detector in (F.7.41).


\medskip\noindent\textbf{An exact normalized detector and its naturality equation.}
The first correspondence in (F.7.35) gives a particularly
short formula for the orbit-twenty detector.  All maps
in this paragraph are the based maps already fixed in
(F.7.19)--(F.7.40).  Thus $x=d^2$, $y=e^2$ is the
specified free basis of $K$, and $p_0:B\to E^+$ is the
row retaining $J_0=\{0,1,2\}$.

For $\gamma=(1,f)\in GT_0^1$, let $j\in\{0,\ldots,4\}$
be the actual finite correction in (F.7.27), and put
\[
 \begin{split}
 W_f&=p_0(d^{-j}f^{-1}e^j)\in K,\\
 A_f&=\ell_x(W_f),\qquad B_f=\ell_y(W_f),\\
 F_0(f)&=y^{B_f}W_f^{-1}x^{A_f}\in K'.
 \end{split}                                                   \tag{F.7.42}
\]
Here $\ell_x,\ell_y:K\to\widehat\Z$ are the
coordinates of the displayed abelianization.
The source of $p_0$ in this formula is legitimate:
$d^{-j}f^{-1}e^j=h_1$ lies in $H\leq B$ by
(F.7.27)--(F.7.28).
The exact normalized representative of the detector is
\[
 \mathcal D_{J_0}(\gamma)
       =\bigl[(x\mapsto x,\quad
                   y\mapsto F_0(f)^{-1}yF_0(f))\bigr]
                    \quad\hbox{in }\mathcal P_K^1.              \tag{F.7.43}
\]
In particular this is a formula in the actual target
pure peripheral group; membership of its normalized
coordinate in $GT_0^1$ is not used.

To prove the formula, use the two actual $H$ cusp
words $c_0=d^5$ and $c_1=e^5=b c_0b^{-1}$.
The row in (F.7.35) and the identification (F.7.36)
give
\[
       p_0(c_0)=(aba)b=x,\qquad
       p_0(c_1)=bxb^{-1}=y .
\]
These images are the previously proved full basis,
so (F.7.28) shows that the actual descended
automorphism is
$\psi_f(x)=x$, $\psi_f(y)=W_fyW_f^{-1}$.
This agrees with $p_0\eta_f^H\iota_0|_K$ because
$p_0\iota_0|_K=1_K$ and $\psi_fp_0=p_0\eta_f^H$
on $H$.  Composing $\psi_f$ with
$\operatorname{Inn}(x^{-A_f})$ gives the representative
in (F.7.43), since
$F_0(f)^{-1}yF_0(f)=
x^{-A_f}W_fyW_f^{-1}x^{A_f}$.
The abelianization of $F_0(f)$ is zero.
Uniqueness is the normalized-pair uniqueness:
two representatives fixing $x$ differ by conjugation
in $\overline{\langle x\rangle}$, and equality of the
images of $y$ then leaves a factor in
$\overline{\langle y\rangle}$; the two zero
abelian coordinates remove both factors.

Consequently the kernel and the equality comparison
with the original normalized action have the exact
descriptions
\[
 \begin{split}
 \ker\mathcal D_{J_0}
   &=\{\,(1,f)\in GT_0^1:
                     W_f=x^{A_f}y^{B_f}\,\},\\
 \{\gamma\in GT_0^1:\mathcal D_{J_0}(\gamma)=\gamma
                                      \text{ in }\mathcal P_K^1\}
   &=\{\,(1,f)\in GT_0^1:
                     y^{B_f}W_f^{-1}x^{A_f}=f\,\}.
 \end{split}                                                   \tag{F.7.44}
\]
For the first equality one can also argue directly.
An inner automorphism fixing $x$ is conjugation
by $x^t$.  Its equality with $\psi_f$ is exactly
$x^{-t}W_f\in C_K(y)=\overline{\langle y\rangle}$,
which forces $t=A_f$ and the remaining exponent
to be $B_f$.
No assertion that the second set in (F.7.44)
is all of $GT_0^1$ is made.

The normalized formula is independent of the lift
of the finite correction.  Replacing $j$ by
$j+5m$, with $m\in\widehat\Z$, gives the exact
identities
\[
 h_1\longmapsto c_0^{-m}h_1c_1^m,\qquad
 W_f\longmapsto x^{-m}W_fy^m,\qquad
 (A_f,B_f)\longmapsto(A_f-m,B_f+m).
                                                               \tag{F.7.45}
\]
Substitution in (F.7.42) leaves $F_0(f)$ unchanged.
This also proves explicitly that this normalization
absorbs the permissible $H$-inner change.

For completeness the exact naturality defect can
be written before any filling.  Put
$\eta_j=\operatorname{Inn}(d^{-j})\eta_f$ and
\[
 L_f=x f^{-1}b,\qquad
 T_f=L_f f^{-1}a L_f^{-1}.
\]
Since $z=xb$, these are $\eta_f(z)$ and
$\eta_f(\tau)$, respectively.  The actual two
basis values are
\[
 \begin{split}
 \eta_j\iota_0(x)
     &=d^{-j}(T_f f^{-1}b)^2d^j,\\
 \eta_j\iota_0(y)
     &=d^{-j}(f^{-1}bT_f)^2d^j .
 \end{split}                                                   \tag{F.7.46}
\]
Define two elements of $H$ by
\[
 \begin{split}
 \Delta_x(f)&=\eta_j\iota_0(x)\,\iota_0(x)^{-1},\\
 \Delta_y(f)&=\eta_j\iota_0(y)
              \bigl(\iota_0(f)^{-1}
                       \iota_0(y)\iota_0(f)\bigr)^{-1}.
 \end{split}
\]
For every one of the four exact filling maps, the
actual representative of (F.7.40) satisfies
\[
 \begin{split}
 (p_i\eta_f^H\iota_0)(x)&=p_i(\Delta_x(f))\,x,\\
 (p_i\eta_f^H\iota_0)(y)&=
                   p_i(\Delta_y(f))\,f^{-1}yf .
 \end{split}                                                   \tag{F.7.47}
\]
For $i=0$ the first defect is killed, by the
$c_0$ computation above.  The second comparison,
including its permitted inner normalization, is
exactly (F.7.44).  A section of the quotient map
supplies the evaluations in (F.7.47); it does not
supply an additional intertwining identity.

There is also an explicit finite certificate for
the size of that section's image.  The same
$\iota_0$ maps $K$ into $K$, with exact word values
\[
 \begin{split}
 \iota_0(x)&=xyx^{-1}y^{-1}x^{-1}yx,\\
 \iota_0(y)&=yx^{-1}y^{-1}xy^{-1}x^{-1}.
 \end{split}                                                   \tag{F.7.48}
\]
Reduction in $C_2*C_3$ verifies both identities.
Their abelianization matrix is
$\left(\begin{smallmatrix}0&-1\\1&-1\end{smallmatrix}\right)$,
so the induced endomorphism of each free pro-$p$
quotient is an automorphism by its Frattini
criterion and the Hopf property.  Nevertheless
$q_5\iota_0(K)=1$, whereas $q_5(K)=A_5$:
the images of $x=d^2$ and $y=e^2$ generate $A_5$,
as the finite certificate verifies.
Thus $\iota_0|_K$ is an injective, proper full
profinite self-embedding with the displayed
finite quotient obstruction to surjectivity.
Its common-section property is fully compatible
with this exact comparison.

The executable
\path{code/part_a/modular_detector_defect_f7.py}
checks the two based cusp images, all four common
section basis images, the words (F.7.48), their
$S_3$ and $A_5$ images, and the normalization,
kernel, and lift-independence identities on all
$4\,373$ reduced words of length at most seven.
These word checks verify the displayed formulas;
they make no faithfulness assertion about
(F.7.43) or the full five-detector family.


\medskip\noindent\textbf{One cone conjugator and all five exact filling equations.}
The kernel equation (F.7.44) has an exact formulation
using a single element of $H$.  The full profinite
generation assertion needed in this formulation
follows from the actual degree-one cusp of a finite
orbifold cover.

First, every closed subgroup $L\leq E$ satisfying
\[
              d\in L,\qquad d^6\in\overline{[L,L]}
                                                               \tag{F.7.49}
\]
is $E$.  Indeed a proper such $L$ is contained in a
proper open subgroup $U$: separate a point outside
$L$ in a finite quotient and take the inverse image
of the image of $L$.  The subgroup
$U\cap\Delta$ has the same finite index, and its full
completion is $U$.  In the corresponding connected
orbifold cover, $d\in U\cap\Delta$ gives a width-one
cusp through the base point.  A proper cover has
at least one more cusp, since the remaining cosets
cannot lie in that length-one cycle of $d$.
Its cone orders divide two or three.
For any prime $p>3$, the cone classes vanish in
$H_1(U,\F_p)$; the cusp classes have precisely their
total-sum relation, together with independent
handle classes.  Each cusp class is therefore
nonzero.  The class of $d^6$ is six times the
specified width-one cusp and is nonzero, contradicting
(F.7.49).

It follows that every continuous endomorphism
$\rho:E\to E$ satisfying
$\rho(d)=t d^\mu t^{-1}$ for a full unit
$\mu\in\widehat\Z^\times$ is an automorphism.
The compact image $L=t^{-1}\rho(E)t$ contains
$d^\mu$ and hence its $\mu^{-1}$ power $d$.
Since $E^{\rm ab}=C_2\times C_3$, the element
$d^6$ belongs to the closed commutator subgroup
of $E$.  Thus $d^{6\mu}\in\overline{[L,L]}$,
and its $\mu^{-1}$ power gives (F.7.49).
Surjectivity follows, and finite generation of $E$
gives Hopficity.  Both the conjugation and the
unit-root operation here are performed on the
specified image subgroup; no pro-$p$ quotient
replaces $E$.

Define the following compact solution space:
\[
 \mathscr S=
 \{\,g\in H:(d g b^{-1}g^{-1})^2=1\,\},\qquad
 \eta_g(d)=d,\quad
 \eta_g(b)=g b g^{-1},\quad
 \eta_g(a)=d g b^{-1}g^{-1}.
                                                               \tag{F.7.50}
\]
The equation gives the $a^2$ relation, and the
$b^3$ relation holds by conjugation.  Hence the
displayed substitutions define an endomorphism
of $E$, and the preceding full boundary test makes
it an automorphism.  Its action on $A_5$ is the
identity because $g\in H$.  Consequently its outer
class is the element of $GT_0^1$ specified by
(F.7.23).

Every element of that $GT_0^1$ has this form.
For $\gamma=(1,f)$, use the actual
$\eta_j=\operatorname{Inn}(d^{-j})\eta_f$ in
(F.7.27).  The full free profinite product torsion
theorem makes $\eta_j(b)$ conjugate to $b$ or
$b^{-1}$; these are the only order-three factor
elements in $C_2\amalg C_3$
\cite[Theorem A and Section 3]{A-HerfortRibes1985}.
Fixing $d$ forces the identity on
$E^{\rm ab}=C_6$, which excludes $b^{-1}$.
Choose $u$ with $\eta_j(b)=ubu^{-1}$.
Since the quotient action is the identity,
$q_5(u)\in C_{A_5}(q_5(b))=\langle q_5(b)\rangle$.
Write $q_5(u)=q_5(b)^r$ and replace $u$ by
$g_f=ub^{-r}\in H$.  The same source's full
free-factor centralizer theorem gives
$C_E(b)=\langle b\rangle$
\cite[Theorem B and Section 4]{A-HerfortRibes1985}.
As $H\cap\langle b\rangle=1$, this $g_f$ is unique.
The conjugation map $H\to E$, $g\mapsto gbg^{-1}$,
is an embedding of a compact space and thus has
a continuous inverse on its image.
We have the exact identities
\[
 \eta_j=\eta_{g_f},\qquad
 h_1=g_f b g_f^{-1}b^{-1},\qquad
 f=b d^j g_f b^{-1}g_f^{-1}d^{-j}.
                                                               \tag{F.7.51}
\]
The last formula retains the particular finite
$j$ selected for the original normalized pair;
it follows from $\eta_f=\operatorname{Inn}(d^j)\eta_j$.

Two solutions $g,h\in\mathscr S$ have the same
outer class exactly when
\[
                   h=d^{5t}g\quad(t\in\widehat\Z).
 \qquad
 \mathscr S/\bigl(g\sim d^{5t}g\bigr)
                         \xrightarrow{\ \sim\ }GT_0^1.
                                                               \tag{F.7.52}
\]
To prove the assertion, an inner comparison of
$\eta_g,\eta_h$ must centralize $d$, so its
conjugator lies in $\overline{\langle d\rangle}$
by (F.7.24).  Both quotient actions are the
identity; since $Z(A_5)=1$, that conjugator is
$d^{5t}\in H$.  Comparing the two images of $b$
then gives $h=d^{5t}g$, because the possible
remaining factor lies in
$H\cap C_E(b)=1$.  Conversely left multiplication
by $d^{5t}$ changes $\eta_g$ by precisely
$\operatorname{Inn}(d^{5t})$ and preserves
(F.7.50).
This proves the bijection, and compactness gives
the stated homeomorphism of the solution-orbit
space with the compact group on the right.

For the first detector the exponent condition
sharpens to one procyclic condition:
\[
 \begin{split}
 \ker\mathcal D_{J_0}
   &=\{\, (1,f):W_f=x^{A_f}y^{-A_f}\,\}\\
   &=\{\, (1,f):p_0(g_f)\in
                            \overline{\langle x\rangle}\,\},\\
 \mathscr S\cap\ker(p_0|_H)
   &\xrightarrow{\ \sim\ }\ker\mathcal D_{J_0}.
 \end{split}                                                   \tag{F.7.53}
\]
For the proof let $\varphi_g$ be the actual
automorphism of $E^+$ descended from $\eta_g$
through $p_0$.  It sends $b$ to
$p_0(g)b p_0(g)^{-1}$, and restricts to the
surface automorphism in (F.7.43).
When its outer restriction is the identity,
that restriction is $\operatorname{Inn}(x^{A_f})$,
by (F.7.44).  Since $C_{E^+}(K)=1$, two
automorphisms of $E^+$ with this same restriction
are equal: their quotient fixes $K$ pointwise,
and its difference on each element centralizes
$K$.  It follows that
$\varphi_g=\operatorname{Inn}(x^{A_f})$.
Comparing the images of $b$ gives
$p_0(g)=x^{A_f}$, since
$K\cap C_{E^+}(b)=1$.
Therefore
$W_f=x^{A_f}b x^{-A_f}b^{-1}=x^{A_f}y^{-A_f}$.
Conversely this expression is exactly the inner
kernel equation (F.7.44).
Finally (F.7.52) replaces a solution with
$p_0(g)=x^t$ uniquely by $d^{-5t}g$ with
$p_0(d^{-5t}g)=1$.  This gives the last
homeomorphism in (F.7.53).
The group law on this slice is explicitly
$g\star h=\eta_g(h)g$: here $\eta_g$ descends
to the identity through $p_0$, so the product
still lies in the slice.

The other four representative fillings can be
written in the same actual coordinates.
An independent index-sixty Schreier calculation
gives $H$ with $61$ initial generators and
$120$ rewritten relators.  Fifty single-occurrence
Tietze eliminations leave eleven generators and
no relators.  For each of the five triples in
(F.7.32), the nine omitted cusp relators give
nine more such eliminations, leaving a full
free basis $P,Q$ with the following values:
\[
 \begin{array}{c|c|ccc|cc}
 J&m&c_0&c_m&c_{\rm third}&P\text{ in }x,y&Q\text{ in }x,y\\ \hline
 \{0,1,2\}&1&QP&Q^{-1}&P^{-1}&yx&y^{-1}\\
 \{0,1,4\}&1&Q&P^{-1}&Q^{-1}P&y^{-1}&x\\
 \{0,1,8\}&1&QP&Q^{-1}&P^{-1}&yx&y^{-1}\\
 \{0,1,10\}&1&Q&P^{-1}&Q^{-1}P&y^{-1}&x\\
 \{0,4,9\}&4&P&P^{-1}Q&Q^{-1}&x&xy
 \end{array}                                                   \tag{F.7.54}
\]
The last two columns are inverse substitutions:
they identify $c_0,c_m$ with the specified full
basis $x,y$ of $K$.
All omitted cusp words go to one.
Thus the resulting maps
$\pi_J:H\to K$ are precisely the actual fillings
(F.7.29) followed by these based quotient
isomorphisms.  The complete original generators,
their words in $a,b$, all relators, all
eliminations, and both inverse basis maps
are recorded in the certificate.
For $J_0$, this is exactly $p_0|_H$,
as both identifications send $c_0,c_1$ to $x,y$.

For $g\in\mathscr S$ put $h(g)=g b g^{-1}b^{-1}\in H$.
Use the original $y=e^2\in E$ in the following
source word:
\[
 \begin{array}{c|c}
 J& V_J(g)\\ \hline
 \{0,1,2\},\{0,1,4\},\{0,1,8\},\{0,1,10\}
       &\pi_J(h(g))\\
 \{0,4,9\}&
          \pi_J\bigl(h(g)y h(g)^{-1}y^{-1}\bigr).
 \end{array}                                                   \tag{F.7.55}
\]
The whole word in the last row lies in $H$,
because $H$ is normal; the map $\pi_J$ is
applied to that word and is not evaluated
separately on $y$ outside its source.
To check the row, the exact label-four
whisker is $(ba)^2=y$, so $c_4=yc_0y^{-1}$.
The substitutions (F.7.50) fix $d$ and $x$,
and give $\eta_g(y)=h(g)y h(g)^{-1}$.
Consequently the conjugator for $c_4$ is
$h(g)y h(g)^{-1}y^{-1}$.
For the first four rows the conjugator for
$c_1=e^5$ is $h(g)$, as in (F.7.51).
Every descended map therefore fixes $x$ and
sends $y$ to $V_J(g)yV_J(g)^{-1}$.
If $A_J=\ell_x(V_J(g))$ and
$B_J=\ell_y(V_J(g))$, all five exact kernel
equations are
\[
 \begin{split}
 \mathcal D_J([\eta_g])&=
       [\,x\mapsto x,\ 
          y\mapsto F_J(g)^{-1}yF_J(g)\,],\\
 F_J(g)&=y^{B_J}V_J(g)^{-1}x^{A_J},\\
 \mathcal D_J([\eta_g])=1
             &\quad\Longleftrightarrow\quad
                   V_J(g)=x^{A_J}y^{B_J}.
 \end{split}                                                   \tag{F.7.56}
\]
Here $[\eta_g]$ is read as its exact $GT_0^1$
class under (F.7.23).
This supplies the same single full profinite
$g$ in every row, together with the
quadratic equation (F.7.50).

The four $C_3$-stable triples of (F.7.34) admit
one further common word.  With the same
$\tau$ from (F.7.33), define
\[
 R(g)=g^{-1}\eta_g(\tau)g\tau^{-1}\in H.
 \qquad
 \mathcal D_{J_i}([\eta_g])=1
       \quad\Longleftrightarrow\quad p_i(R(g))=1
       \quad(0\leq i\leq3).
                                                               \tag{F.7.57}
\]
Indeed $r=b$ is conjugated by $g$, while
$s=\tau b\tau^{-1}$ is conjugated by
$v(g)=\eta_g(\tau)g\tau^{-1}$.
After filling, $r,s$ generate $E^+$.
An inner surface action extends uniquely across
$E^+$ and its conjugator must be $p_i(g)$,
by $C_{E^+}(K)=1$ and
$K\cap C_{E^+}(b)=1$.
Its agreement on $s$ is then exactly
$p_i(v(g))=p_i(g)$, because the other finite
factor also has trivial centralizer intersection
with $K$.  This is (F.7.57), in both directions.
Replacing $g$ by $d^{5t}g$ and $\eta_g$ by its
specified inner change leaves $R(g)$ unchanged,
so this common-word comparison retains the
outer normalization.

The actual geometric residual group in (F.7.25)
therefore maps into the simultaneous solution
set (F.7.50), (F.7.53), and (F.7.56), with the
common-word constraints (F.7.57).
Triviality of that solution set has not been
proved.  The nontrivial common source kernel
in (F.7.41) supplies no solution to the additional
quadratic equation.  Conversely the exact
section and the primewise images in (F.7.48)
supply no step replacing these full quotient
equations.

The executed standard-library program
\path{code/part_a/modular_five_fillings_f7.py}
records all $95$ Tietze eliminations, all five
inverse full-basis maps, every original cusp
word, the order-three finite centralizer,
the five $j$-recovery word identities in
(F.7.51), and the exact label-four conjugator.
The orbifold generation test, the full torsion
and centralizer inputs, and the simultaneous
solution-space identifications are proved above.


This audit concerns the comparison requested for
\S30, folios 194--196, scans 383--387. F.6's
arithmetic open-image proof applies to $A_5$ without
using (F.7.1), so that proved arithmetic assertion
does not depend on the unresolved intersection.
\end{proof}
